\chapter{Conclusion and Future Work}

\section{Overview}

This project set out to build a system that does what existing AI developer
tools do not: close the loop between detecting a software failure and
deploying a verified fix for it, with a human involved as a single decision
point rather than as the pipeline itself. The delivered platform, Nexis,
realises that goal as a working, verified system rather than a design
proposal.

The central deliverable is the nine-agent closed-loop self-healing pipeline,
durably orchestrated by Temporal \cite{temporal}. A five-agent Execution Team
(Architect, Backend, QA, DevOps, Data Engineer) performs the engineering
work, while a four-agent Self-Healing Loop (Sentinel, Pathfinder,
Synthesiser, Validator) detects faults, localises root causes over the Neo4j
dependency graph \cite{neo4j}, sandbox-validates every candidate patch, and
routes the result through a severity-classified Approval Gate offering
Approve, Reject, and Modify decisions. This pipeline was not merely
implemented but exercised live: Chapter 6 reports two complete end-to-end
recovery runs against a real local language model \cite{ollama} -- one in
which the 120-second approval window genuinely elapsed and the run was
timeout-rejected exactly as the safety design intends, and one in which a
detected fault proceeded through planning, patch generation, property-based
test validation, human approval, and deployment to a Succeeded terminal
state in approximately one minute and forty seconds.

The second major deliverable is the from-scratch Docker preview-deploy
engine \cite{docker}. Given a connected repository, it clones the code,
detects the technology stack, generates a working Dockerfile when none
exists, builds and runs the container under resource limits, and returns a
live preview URL; build and health-check failures feed back into the same
self-healing pipeline as any other incident source. This too was proven
live: a real public repository with no Dockerfile was cloned, containerised,
built, served, HTTP-verified, and cleanly stopped in a single recorded run.

Around these two engines sits a complete multi-tenant platform: JWT and SSO
authentication with MFA and session management \cite{workos}, role-based
access control enforced down to the database row via PostgreSQL Row-Level
Security \cite{pgrls}, the heuristic role-recommendation engine that turns
audit-log evidence into explainable privilege suggestions, Stripe-based
billing \cite{stripe}, webhook-driven integrations with GitHub, Sentry,
Datadog, PagerDuty, and Slack \cite{sentry,datadog,pagerduty,slackapi}, and
an observability surface (dashboards, cost tracking, system health probes,
hash-chained audit log) built on OpenTelemetry \cite{otel}. All Go modules,
the Python causal-inference service, and the Next.js console pass their full
build, vet, lint, and test suites, as documented in Chapter 6.

\section{Limitations of the Project}

Academic integrity requires that the boundary between what was demonstrated
and what was not be stated plainly. The following limitations are real,
known, and deliberate; none is hidden behind softer wording.

\begin{itemize}
  \item \textbf{No production credentials were exercised.} No real GitHub
  App or OpenAI production credentials existed in the evaluation environment
  used for this report. The live evidence was produced with local Ollama
  models \cite{ollama} and an unauthenticated clone of a public repository;
  a fully credentialed production path -- ending in a real merged GitHub
  pull request -- was not exercised end-to-end.
  \item \textbf{The platform is not yet live on its domain.} The
  \texttt{nexis.sh} domain is registered but the platform is not yet
  deployed to it. All evidence in this report comes from local, native
  execution. Docker was unavailable in the build environment, so Podman
  \cite{podman} was substituted as a Docker-API-compatible runtime and shown
  to behave identically.
  \item \textbf{Fine-tuning is a manual step by design.} The RLHF
  data-export pipeline is real and tested -- every terminal approval
  decision, including engineer-edited diffs, is persisted and exportable as
  JSONL -- but actually submitting a fine-tuning job against that export is
  a deliberate manual step involving real money and a real external API
  call, not an automated one.
  \item \textbf{One of twenty-seven demo scenarios is fully wired.} The
  Live-Demo console lists 27 curated fault scenarios, of which exactly one
  is wired end-to-end to a real fixture and classification path today; the
  remaining 26 are UI-complete but backend-pending.
  \item \textbf{The causal ranking is a proxy, not a fitted causal model.}
  The causal-inference sidecar ranks root-cause candidates using an honest
  statistical and graph-structural scoring method with a transparent
  per-component confidence breakdown. It is explicitly not a fitted DoWhy
  structural causal model \cite{dowhy}, because no interventional production
  data exists against which such a model could be fitted; claiming otherwise
  would be an overstatement.
  \item \textbf{A known unrelated bug was left unfixed.} A pre-existing bug
  in the workspace-onboarding progress poller (a 404 from
  \texttt{/v1/workspaces/\{id\}/events}) was found during evaluation and
  intentionally left unfixed as out of scope; the workspace itself
  provisions successfully server-side regardless.
  \item \textbf{No formal load testing was performed.} As noted in
  Chapter 6, no formal concurrent-load performance benchmark was run; the
  reported latencies come from single-run live executions, not from a
  controlled load harness.
\end{itemize}

\section{Future Work}

Each item below is grounded in the existing architecture and represents a
concrete next step rather than a generic aspiration.

\begin{itemize}
  \item \textbf{Credentialed production deployment.} Deploy the platform to
  the registered \texttt{nexis.sh} domain behind a real GitHub App
  installation and a paid OpenAI account, so that the full production path
  -- authenticated clone, real pull-request creation, real merge-triggered
  redeploy -- is exercised with production credentials end-to-end.
  \item \textbf{Complete the scenario library.} Wire the remaining 26
  Live-Demo scenario cards to real fixtures and classification paths,
  extending the existing 11-scenario fault-injection library across the
  application, database, deploy, observability, and security categories it
  already enumerates.
  \item \textbf{Finer-grained Sentinel baselines.} The Sentinel's
  statistical spike rule currently maintains its EWMA mean and variance
  baseline per organisation; extending it to per-(organisation, service)
  baselines would let a noisy service stop masking a quiet one.
  \item \textbf{A fitted causal model.} Once production incident data with
  genuine interventions accumulates, fit a real DoWhy structural causal
  model \cite{dowhy} and evaluate it against the current graph-evidence
  proxy, which was designed from the outset to be replaceable.
  \item \textbf{Execute a fine-tuning run.} Once the accumulated RLHF export
  contains enough approve/reject/modify examples, submit an actual
  fine-tuning job against it and measure whether the tuned model reduces
  rejection and modification rates at the approval gate.
  \item \textbf{Real preview subdomains.} Add a reverse-proxy layer (for
  example Traefik) in front of the deploy engine so that preview
  deployments receive stable per-project subdomains instead of
  \texttt{localhost} ports, which is a prerequisite for sharing preview
  URLs beyond the local machine.
  \item \textbf{Formal load benchmarking.} Run a controlled concurrent-load
  benchmark against the control plane and the recovery pipeline to
  characterise throughput, queue behaviour under simultaneous incidents,
  and Temporal worker saturation.
  \item \textbf{Human-subjects usability study.} Conduct the planned
  NASA-TLX cognitive-workload study \cite{nasatlx} with external
  participants; the survey instrument is already integrated into the
  post-incident flow, so the remaining work is recruitment and analysis
  rather than implementation.
\end{itemize}

\section{Conclusion}

The central claim of this report is deliberately narrow and deliberately
strong: that a software fault can travel from automatic detection, through
multi-agent diagnosis, patch generation, sandboxed validation, and a single
human approval decision, to a deployed fix -- inside one continuous,
durable, observable pipeline. That claim is not asserted; it is demonstrated
with reproducible, live evidence. One recorded run shows the loop refusing
to proceed when its human gate was not answered in time, and another shows
the same loop carrying a real fault to an approved, deployed resolution in
under two minutes, with every agent step, token count, and decision recorded
in a durable timeline. The preview-deploy engine extends the same principle
to infrastructure, turning its own build failures into self-healing targets.
Equally important is what this report declines to claim. The honest
accounting in Section 7.2 -- local models rather than production
credentials, a registered but not yet live domain, a proxy rather than a
fitted causal model, one wired demo scenario out of twenty-seven -- is not a
list of apologies but part of the project's rigor: the boundary of the
evidence is drawn exactly where the evidence stops. Within that boundary,
Nexis stands as a complete, working demonstration that closed-loop fault
recovery, gated by a single human decision point, is buildable today.
